\documentclass{article}
\usepackage[utf8]{inputenc}
\usepackage[spanish]{babel}
\usepackage{listings}
\usepackage{xcolor}
\usepackage{geometry}
\geometry{a4paper, margin=1in}

\% Configuración para código C++
\definecolor{codegreen}{rgb}{0,0.6,0}        
\definecolor{codegray}{rgb}{0.5,0.5,0.5}     
\definecolor{codepurple}{rgb}{0.58,0,0.82}   
\definecolor{backcolour}{rgb}{0.95,0.95,0.92}

\lstset{
    backgroundcolor=\color{backcolour},      
    commentstyle=\color{codegreen},
    keywordstyle=\color{magenta},
    numberstyle=\tiny\color{codegray},       
    stringstyle=\color{codepurple},
    basicstyle=\ttfamily\small,
    breakatwhitespace=false,
    breaklines=true,
    captionpos=b,
    keepspaces=true,
    numbers=left,
    numbersep=5pt,
    showspaces=false,
    showstringspaces=false,
    showtabs=false,
    tabsize=2,
    language=C++,
    inputencoding=utf8,
    extendedchars=true,
    literate={á}{{\'a}}1 {é}{{\'e}}1 {í}{{\'i}}1 {ó}{{\'o}}1 {ú}{{\'u}}1 {ñ}{{\~n}}1 {¿}{{\textquestiondown}}1 {¡}{{\textexclamdown}}1
}

\% Macros para las secciones de los apuntes
\newcommand{\quees}[1]{
    \section*{🤔 ¿QUÉ ES?}
    \hrulefill \\
    \#1
}

\newcommand{\dichodeotraforma}[1]{
    \section*{💬 DICHO DE OTRA FORMA}
    \hrulefill \\
    \#1
}

\newcommand{\diagrama}[1]{
    \section*{🎨 DIAGRAMA}
    \hrulefill \\
    \#1
}

\newcommand{\comofunciona}[1]{
    \section*{⚙️ ¿CÓMO FUNCIONA?}
    \hrulefill \\
    \#1
}

\newcommand{\complejidad}[1]{
    \section*{📱 COMPLEJIDAD (Big-O)}
    \hrulefill \\
    \#1
}

\newcommand{\entrevistas}[1]{
    \section*{🎯 EN ENTREVISTAS}
    \hrulefill \\
    \#1
}

\newcommand{\resumen}[1]{
    \section*{📋 RESUMEN RÁPIDO}
    \hrulefill \\
    \#1
}

\begin{document}

\title{Quick Sort (Ordenamiento Rápido)}
\author{Aldo Zepeda Montoya}
\date{\today}
\maketitle

\subsection*{CATEGORÍA: 04\_Algoritmos\_Ordenamiento}

\quees{
Elegir a un jugador de referencia ⚽ en el equipo (el Pivot). 

Una vez que eliges al Pivot (por ejemplo, el último de la fila), 
separas a todos los demás en dos grupos:
1. Los que son "peores" (número menor) van a la izquierda del Pivot.
2. Los que son "mejores" (número mayor) van a la derecha del Pivot.

Ahora el Pivot quedó en su lugar correcto para siempre. Luego, haces
lo mismo con los dos grupos que quedaron a los lados, hasta que todos
estén ordenados.
}

\dichodeotraforma{
Quick Sort es un algoritmo de ordenamiento por comparación, 
recursivo e In-place. 

Se basa en la operación "Partition": elegir un elemento como Pivot y 
reorganizar el array de modo que los menores queden a la izquierda y 
los mayores a la derecha. 

Aunque su peor caso es O(n^2), en la práctica es mucho más rápido que 
Merge Sort debido a que tiene una excelente "Cache Locality" y no 
requiere memoria extra significativa.
}

\diagrama{\begin{verbatim}
Array: [3, 6, 8, 10, 1, 2, 5]  (Pivot = 5)

Partition:
- Elementos < 5: [3, 1, 2]
- Pivot: 5
- Elementos > 5: [6, 8, 10]

Estado: [3, 1, 2] [5] [6, 8, 10]
(El 5 ya no se mueve más).

Se repite el proceso para [3, 1, 2] y para [6, 8, 10].
\end{verbatim}}

\begin{lstlisting}[language=C++]
1. Caso base: Si el rango tiene 0 o 1 elementos, ya está ordenado.
2. Partition:
   - Eliges un Pivot (el primero, el último, el medio o uno aleatorio).
   - Reacomodas el array: elementos < Pivot a la izquierda, > a la der.
   - Devuelves la posición final del Pivot.
3. Llamas recursivamente a Quick Sort para la mitad izquierda.
4. Llamas recursivamente a Quick Sort para la mitad derecha.

📝 PSEUDOCÓDIGO
──────────────────────────────────────────────────────────────

función quickSort(arr, inicio, fin):
    SI inicio < fin:
        p_idx = partition(arr, inicio, fin)
        quickSort(arr, inicio, p_idx - 1)
        quickSort(arr, p_idx + 1, fin)

función partition(arr, inicio, fin):
    pivot = arr[fin]
    i = inicio - 1
    PARA j desde inicio hasta fin - 1:
        SI arr[j] < pivot:
            i++
            intercambiar(arr[i], arr[j])
    intercambiar(arr[i+1], arr[fin])
    return i + 1

💻 CÓDIGO C++
──────────────────────────────────────────────────────────────

#include <iostream>
#include <vector>
using namespace std;

int partition(vector<int>& arr, int low, int high) {
    int pivot = arr[high]; // Elegimos el último como pivot
    int i = low - 1;       // Índice del elemento más pequeño

    for (int j = low; j < high; j++) {
        if (arr[j] < pivot) {
            i++;
            swap(arr[i], arr[j]);
        }
    }
    swap(arr[i + 1], arr[high]);
    return i + 1;
}

void quickSort(vector<int>& arr, int low, int high) {
    if (low < high) {
        int pi = partition(arr, low, high);
        quickSort(arr, low, pi - 1);
        quickSort(arr, pi + 1, high);
    }
}

int main() {
    vector<int> data = {10, 7, 8, 9, 1, 5};
    quickSort(data, 0, data.size() - 1);
    for(int x : data) cout << x << " "; // 1 5 7 8 9 10
    return 0;
}

⏱️ COMPLEJIDAD (Big-O)
──────────────────────────────────────────────────────────────

  Caso        │ Tiempo        │ Cuándo ocurre
  ────────────┼───────────────┼────────────────────────────────
  Best        │ O(n log n)    │ El pivot siempre parte a la mitad
  Average     │ O(n log n)    │ Pivot aleatorio / balanceado
  Worst       │ O(n^2)         │ Array ya ordenado + pivot extremo
  Espacio     │ O(log n)      │ Stack de recursión
\end{lstlisting}

\begin{lstlisting}[language=C++]
✅ std::sort: En C++, `std::sort` usa "Introsort", que empieza con 
   Quick Sort y cambia a Heap Sort si la recursión es muy profunda.

💡 Pivot Selection: Para evitar el peor caso O(n^2), menciona que se 
   puede usar un Pivot aleatorio o la técnica "Median of Three".

⚠️ No es estable: A diferencia de Merge Sort, Quick Sort NO garantiza
   mantener el orden relativo de elementos iguales.

🔑 In-place: Su gran ventaja es que no necesita arrays extras, solo 
   memoria para las llamadas recursivas (el stack).
\end{lstlisting}

\resumen{
• El más rápido en la práctica.
• Usa la estrategia "Divide y Vencerás".
• In-place (ahorra memoria).
• Peor caso O(n^2) pero muy raro con buen pivot.
• No es estable.
• Operación clave: `partition`.
════════════════════════════════════════════════════════════════
}

\end{document}